\chapter{Communication Complexity — Helper}
%%% generated by the blueprint pipeline from TCSlib.CommunicationComplexity.DeterministicCC.Helper
\section{Overview}

This module provides Boolean helper definitions used throughout the communication
complexity library: the type of $n$-bit Boolean inputs, the $\pm 1$ sign encoding
of Boolean values, and coordinate-flipping operations, together with their basic
algebraic identities.

%%%%%%%%%%%%%%%%%%%%%%%%%%%%%%%%%%%%%%%%%%%%%%%%%%%%%%%%%%%%%%%%%%%%%%%%%%%%%%%
\section{Declarations}

\begin{definition}[$n$-bit Boolean input type]
\lean{CommunicationComplexity.BoolInput}
\label{CommunicationComplexity.BoolInput}
\leanok
$\texttt{BoolInput}(n)$ is the type of $n$-bit Boolean inputs, defined as the function
type $\mathrm{Fin}\,n \to \mathrm{Bool}$.
\end{definition}

\begin{definition}[All-zero Boolean input]
\lean{CommunicationComplexity.zeroInput}
\label{CommunicationComplexity.zeroInput}
\leanok
\uses{CommunicationComplexity.BoolInput}
$\texttt{zeroInput}(n) : \texttt{BoolInput}(n)$ is the constant $\mathrm{false}$ function,
i.e.\ the $n$-bit string consisting entirely of zeros.
\end{definition}

\begin{lemma}[Nonzero Boolean input has a true coordinate]
\lean{CommunicationComplexity.exists_true_of_ne_zeroInput}
\label{CommunicationComplexity.exists_true_of_ne_zeroInput}
\leanok
\uses{CommunicationComplexity.BoolInput, CommunicationComplexity.boolSign, CommunicationComplexity.zeroInput}
\difficulty{3}
Let $x$ be an $n$-bit Boolean input, that is, a function assigning a Boolean value to
each of the $n$ coordinates. If $x$ is not the all-zero input — the constant function
taking the value $\mathrm{false}$ everywhere — then there exists a coordinate $i$ at
which $x$ takes the value $\mathrm{true}$.
\end{lemma}

\begin{definition}[Boolean sign]
\lean{CommunicationComplexity.boolSign}
\label{CommunicationComplexity.boolSign}
\leanok
$\texttt{boolSign} : \mathrm{Bool} \to \mathbb{R}$ assigns the real value $1$ to
$\mathrm{false}$ and $-1$ to $\mathrm{true}$, embedding Boolean values into the
multiplicative group $\{\pm 1\} \subset \mathbb{R}$.
\end{definition}

\begin{lemma}[Boolean sign turns XOR into multiplication]
\lean{CommunicationComplexity.boolSign_xor}
\label{CommunicationComplexity.boolSign_xor}
\leanok
\uses{CommunicationComplexity.boolSign}
\difficulty{2}
Write $\mathrm{sgn} : \mathrm{Bool} \to \bbr$ for the Boolean sign map, which sends
$\mathrm{false}$ to $1$ and $\mathrm{true}$ to $-1$. Then for all Boolean values $a, b$,
\[
  \mathrm{sgn}(a \oplus b) \;=\; \mathrm{sgn}(a)\,\mathrm{sgn}(b),
\]
where $a \oplus b$ denotes the exclusive-or of $a$ and $b$; equivalently, $\mathrm{sgn}$
is a homomorphism from $(\mathrm{Bool}, \oplus)$ to the multiplicative group $\{\pm
1\}$.
\end{lemma}

\begin{lemma}[Boolean sign turns parity into a product of signs]
\lean{CommunicationComplexity.boolSign_sum}
\label{CommunicationComplexity.boolSign_sum}
\leanok
\uses{CommunicationComplexity.boolSign, CommunicationComplexity.boolSign_xor}
\difficulty{4}
Let $\alpha$ be a type, let $s$ be a finite set of indices from $\alpha$, and let $f :
\alpha \to \{\mathrm{false}, \mathrm{true}\}$ be a Boolean-valued function. Write
$\sigma$ for the Boolean sign, so that $\sigma(\mathrm{false}) = 1$ and
$\sigma(\mathrm{true}) = -1$, and let $\bigoplus_{i \in s} f(i)$ denote the sum of the
values $f(i)$ taken in the Booleans under exclusive-or, so that this sum equals
$\mathrm{true}$ precisely when an odd number of indices $i \in s$ satisfy $f(i) =
\mathrm{true}$. Then the sign of this parity equals the product of the individual signs:
\[
\sigma\!\left(\bigoplus_{i \in s} f(i)\right) \;=\; \prod_{i \in s}
\sigma\bigl(f(i)\bigr).
\]
\end{lemma}

\begin{lemma}[Product of two Boolean signs]
\lean{CommunicationComplexity.boolSign_mul_boolSign_eq_sub_two_indicator}
\label{CommunicationComplexity.boolSign_mul_boolSign_eq_sub_two_indicator}
\leanok
\uses{CommunicationComplexity.boolSign}
\difficulty{2}
For all Boolean values $a$ and $b$, the product of their signs satisfies
\[
  \mathrm{sign}(a)\cdot\mathrm{sign}(b) \;=\; 1 - 2\,\mathbf{1}[a \ne b],
\]
where $\mathrm{sign}(\cdot)$ sends $\mathrm{false}$ to $1$ and $\mathrm{true}$ to $-1$,
and $\mathbf{1}[a \ne b]$ is $1$ when $a$ and $b$ differ and $0$ when they agree.
Equivalently, this product equals $1$ when the two bits are equal and $-1$ when they
disagree.
\end{lemma}

\begin{definition}[Single-coordinate flip]
\lean{CommunicationComplexity.flipAt}
\label{CommunicationComplexity.flipAt}
\leanok
\uses{CommunicationComplexity.BoolInput}
Given $i : \mathrm{Fin}\,n$ and $x : \texttt{BoolInput}(n)$, $\texttt{flipAt}\,i\,x$ is
the Boolean input obtained from $x$ by negating the $i$-th coordinate and leaving all
other coordinates unchanged.
\end{definition}

\begin{lemma}[Value of the single-coordinate flip at the flipped coordinate]
\lean{CommunicationComplexity.flipAt_apply_same}
\label{CommunicationComplexity.flipAt_apply_same}
\leanok
\uses{CommunicationComplexity.BoolInput, CommunicationComplexity.flipAt}
\difficulty{1}
Let $n \in \bbn$, let $i \in \mathrm{Fin}\,n$ be a coordinate index, and let $x$ be an
$n$-bit Boolean input. The single-coordinate flip of $x$ at $i$ — the input obtained
from $x$ by negating its $i$-th coordinate and leaving all others unchanged — has value
$\neg\,x_i$ at coordinate $i$.
\end{lemma}

\begin{lemma}[Flip leaves other coordinates unchanged]
\lean{CommunicationComplexity.flipAt_apply_ne}
\label{CommunicationComplexity.flipAt_apply_ne}
\leanok
\uses{CommunicationComplexity.BoolInput, CommunicationComplexity.flipAt}
\difficulty{1}
Let $x \colon \mathrm{Fin}\,n \to \mathrm{Bool}$ be an $n$-bit Boolean input, and let
$i, j \in \mathrm{Fin}\,n$ be coordinates with $j \ne i$. Then the input obtained from
$x$ by flipping its $i$-th coordinate agrees with $x$ at coordinate $j$; that is, its
value at $j$ equals $x(j)$.
\end{lemma}

\begin{lemma}[Single-coordinate flip is an involution]
\lean{CommunicationComplexity.flipAt_flipAt}
\label{CommunicationComplexity.flipAt_flipAt}
\leanok
\uses{CommunicationComplexity.BoolInput, CommunicationComplexity.flipAt}
\difficulty{3}
Fix $n \in \bbn$, a coordinate $i \in \mathrm{Fin}\,n$, and an $n$-bit Boolean input $x
\colon \mathrm{Fin}\,n \to \mathrm{Bool}$. Writing $\mathrm{flip}_i$ for the operation
that negates the $i$-th coordinate of an input and leaves all other coordinates
unchanged, applying it twice recovers the original input:
\[
\mathrm{flip}_i\bigl(\mathrm{flip}_i(x)\bigr) = x.
\]
\end{lemma}

\begin{lemma}[Bijectivity of a single-coordinate flip]
\lean{CommunicationComplexity.flipAt_bijective}
\label{CommunicationComplexity.flipAt_bijective}
\leanok
\uses{CommunicationComplexity.BoolInput, CommunicationComplexity.flipAt, CommunicationComplexity.flipAt_flipAt}
\difficulty{2}
Fix $n \in \bbn$ and a coordinate $i \in \mathrm{Fin}\,n$. On the space of $n$-bit
Boolean inputs, functions $\mathrm{Fin}\,n \to \mathrm{Bool}$, consider the
single-coordinate flip that negates the $i$-th bit of its argument and leaves every
other bit unchanged. This map is a bijection.
\end{lemma}
